\chapter{The FTC Bridge: Effective Calculus}\label{ch:effective-calculus}

This chapter supplies the finite certificates consumed by the FTC contract; it
is not a second integral theory or a universal completeness theorem.
\begin{verbatim}
primitive + derivative boxes + mesh -> finite telescope -> endpoint overlap
\end{verbatim}
\lean{ComputableAnalysis.DerivativeBoundFTC}
\lean{ComputableAnalysis.ScheduledIntervalRegularOn}
\leanok

\section{Regression ladder}
Affine functions close in one rectangle. The square is the representative
polynomial derivative-bound schedule: \(\int_0^1x^2=1/3\). Its finite
integral and endpoint evaluators overlap; other routine degrees are obtained
by the generic monomial mechanism when needed.
\lean{ComputableAnalysis.ExactFunction.affineFTCExact}
\lean{ComputableAnalysis.Integral.squareIntegralEffectiveFTC_integral_equiv_one_third}
\leanok

\section{Product and trigonometric transport}
The product witness \(F(x)=x\arctan x\) is the first non-polynomial example.
\lean{ComputableAnalysis.IntegralIdentities.coordinateTimesArctanForwardTwoStageFTC}
\leanok

The tangent-chart square integral is complete. The remaining \(\sin(\pi x)^2\)
theorem is an all-depth geometric transport from nested-radical dyadic boxes
to tangent witnesses, including the reciprocal-\(\pi\) normalization. It is
an explicit finite-certificate target, not a missing classical existence proof.
\lean{ComputableAnalysis.SinPiIntegral.ArctanSinPiConstruction.halfIntegral_equiv_reciprocalPi_of_halfAngle_certificate_family}
\leanok

New theorems add an evaluator, finite range/derivative estimate, stage
schedule, validity proof, and one equivalence edge to the representation
chain; pairwise equivalence of every implementation is unnecessary.
